\section{Condiciones de lectura de las variantes}
\label{sec:condiciones-variantes}

Todas las variantes se comparan con la base de \textbf{la misma semilla} ($2026$), en contraste pareado
por participante ($n = 54$) y con los estadísticos de la sección~\ref{sec:estadistica-variantes}. Los
contrastes en kappa coinciden con los de exactitud y no se repiten. Las comparaciones múltiples se
corrigen con Holm dentro de las familias declaradas de antemano (tabla~\ref{tab:familias}); lo que no
pertenece a ninguna familia se rotula exploratorio y no se corrige.

\begin{table}[H]
\centering
\caption{Familias de contrastes de las variantes y de las semillas adicionales, declaradas antes de calcular los resultados.
Cada familia se corrige con Holm por separado. Las dos familias secundarias son informativas.}
\label{tab:familias}
\small
\begin{tabular}{p{3.6cm}p{8.6cm}c}
\toprule
\textbf{Familia} & \textbf{Contenido} & $m$ \\
\midrule
Entre arquitecturas & Seis pares por protocolo y métrica, con la media de tres semillas (la misma estructura que con la semilla $2026$ sola) & 6 \\
Ablaciones intra-sujeto, \eegnet{} y Conformer & Ocho variantes frente a la base intra-sujeto de la semilla 2026: tres subbandas, dos ventanas, sin aumentación, reflexión de canales y ruido en Fourier & 8 \\
Ablaciones intra-sujeto recortadas, LaBraM y EEGMamba & Reflexión de canales, ruido en Fourier y ventana de \seg{0} a \seg{2} (recorte decidido antes de ejecutar) & 3 \\
Alineamiento euclidiano en \loso{} & \loso{} con alineamiento euclidiano frente a la base \loso{} de la semilla 2026, uno por arquitectura & 4 \\
Reflexión de canales en \loso{} & \loso{} con reflexión de canales frente a la misma base, uno por arquitectura & 4 \\
Calibración, base y alineamiento & Cuatro arquitecturas, dos tamaños de calibración y tres métodos; la familia de los modelos base y la de los modelos con alineamiento se corrigen por separado & 24 y 24 \\
Políticas de normalización & Origen frente a las otras dos políticas, por arquitectura & 8 \\
Alineamiento en la transferencia & Tres evaluaciones por cuatro arquitecturas, variante con alineamiento frente a la base de la semilla 2026 & 12 \\
\midrule
Secundarias & Los $22$ contrastes de las ablaciones intra-sujeto a la vez, y los $8$ de alineamiento y reflexión en \loso{} a la vez & 22 y 8 \\
\bottomrule
\end{tabular}
\end{table}

\subsection{Variantes evaluadas}

Las variantes y su implementación se describen en las secciones~\ref{sec:variantes-banda},
\ref{sec:variantes-aug} y~\ref{sec:variantes-ea}. Para leer los resultados basta saber que la variante
sin aumentación anula las tres transformaciones de la base, que la amplitud $0.1$ del ruido en Fourier es
una elección propia y que el alineamiento es transductivo: el participante de prueba se alinea con sus
propios segmentos sin etiqueta.

\subsection{Cota del efecto del recorte de gradiente}
\label{sec:confusor}

Los datos descartan un efecto medio del cambio de régimen (recorte de gradiente, hardware y código)
superior a entre $0.010$ y $0.020$ de exactitud según la arquitectura, no que el efecto sea nulo
(tabla~\ref{tab:confusor}). En \eegnet{}, el Conformer y LaBraM el contraste de una variante con su base
está confundido con ese cambio (tabla~\ref{tab:regimen}); EEGMamba no cambia de recorte, aunque sí de
hardware y de versión del código. El procedimiento de la cota está en la
sección~\ref{sec:cota-confusor}.

\begin{table}[H]
\centering
\caption{Diferencia por participante entre la semilla $2026$ (sin recorte) y la media de las
semillas $2027$ y $2028$ (con recorte), con intervalo al $95$\,\% de percentiles, y cota resultante.
La cota neta resta la misma diferencia en EEGMamba, que no cambia de recorte. Solo existe donde la base cambia
de régimen entre semillas.}
\label{tab:confusor}
\small
\begin{tabular}{llccc}
\toprule
\textbf{Protocolo} & \textbf{Arq.} & Delta bruto, IC $95$\,\% & Cota bruta & Cota neta \\
\midrule
\loso{} & \eegnet{} & $+0.0030$ $[-0.0042,\ +0.0104]$ & 0.0104 & 0.0100 \\
\loso{} & Conformer & $+0.0072$ $[-0.0043,\ +0.0198]$ & 0.0198 & 0.0178 \\
\loso{} & LaBraM & $-0.0016$ $[-0.0097,\ +0.0063]$ & 0.0097 & 0.0142 \\
Intra-sujeto & LaBraM & $+0.0021$ $[-0.0051,\ +0.0092]$ & 0.0092 & 0.0182 \\
\bottomrule
\end{tabular}
\end{table}

Ninguna de las cuatro diferencias medias se distingue de cero (Wilcoxon, $p$ de $0.37$ a $0.91$). En una
media de tres semillas, con una de ellas sin recorte, la contaminación máxima sería la tercera parte, de
$0.003$ a $0.007$. En EEGMamba, que no cambia de recorte, la misma diferencia llega a $+0.009$ en el
protocolo intra-sujeto ($p$ sin corregir de $0.048$), del orden de las cotas, lo que es compatible con
ruido de semilla y con el cambio de hardware y de código. En el protocolo intra-sujeto de \eegnet{} y el
Conformer, que nunca tuvieron recorte, la cota no es estimable. La regla de lectura, fijada de antemano:
un efecto que no supera la cota se declara \emph{no interpretable}, no \emph{sin efecto}; sin cota
estimable, el resultado se informa como condicionado al régimen.

